\documentclass[../../book-main.tex]{subfiles}
\graphicspath{{\subfix{../..}}}

\begin{document}



\section{Gradient Descent}
\section{Back-Propagation}

The most basic class of neural networks are ``multi-layer perceptrons,'' or functions of the form \(\vec{f} \colon \R^{n} \times \Theta \to \R^{p}\), where \(\Theta\) is the so-called ``parameter space'', and have the form
\begin{align}
    \vec{f}(\vec{x}, \bm{\theta}) 
    &= W^{(m)}\vec{\sigma}^{(m)}(W^{(m-1)}(\cdots \vec{\sigma}^{(1)}(W^{(0)}\vec{x} + \vec{b}^{(0)}) \cdots) + \vec{b}^{(m - 1)}) + \vec{b}^{(m)} \\ 
    \text{where} \quad \bm{\theta}
    &= (W^{(0)}, \vec{b}^{(0)}, W^{(1)}, \vec{b}^{(1)}, W^{(2)}, \vec{b}^{(2)}, \dots, W^{(m)}, \vec{b}^{(m)}) \in \Theta,
\end{align}
and \(\vec{\sigma}^{(1)}, \dots, \vec{\sigma}^{(m)}\) are ``activation functions,'' which we treat as generic differentiable nonlinear functions. Here \(\bm{\theta}\) is bolded because it should not be thought of as a scalar, vector, or matrix, but rather as a collection of matrices and vectors, an object we haven't discussed so far in this class. 

More precisely, if we define the functions \((\vec{z}^{(i)})_{i = 0}^{m}\) and \((\vec{f}^{(i)})_{i = 0}^{m}\) as 
\begin{align}
    \vec{z}^{(0)}(\vec{x}, \bm{\theta})
    &= \vec{f}^{(0)}(\vec{x}, \bm{\theta}) = W^{(0)}\vec{x} + \vec{b}^{(0)} \\ 
    \vec{z}^{(i)}(\vec{x}, \bm{\theta})
    &= \vec{f}^{(i)}(\vec{z}^{(i - 1)}(\vec{x}, \bm{\theta}), \bm{\theta}) = W^{(i)}\vec{\sigma}^{(i)}(\vec{z}^{(i - 1)}) + \vec{b}^{(i)}, \qquad \forall\ i \in \{1, \dots, m\},
\end{align}
then this formulation gives 
\begin{equation}
    \vec{f}(\vec{x}, \bm{\theta}) = \vec{z}^{(m)}(\vec{x}, \bm{\theta}) = \vec{f}^{(m)} \circ \cdots \circ \vec{f}^{(0)}(\vec{x}, \bm{\theta}).
\end{equation}

When we say that we ``train'' this neural network, we mean just that we optimize \(\bm{\theta}\) to minimize some loss function evaluated on the data. Given a (differentiable) loss function \(\ell \colon \R^{p} \times \R^{p} \to \R\) and a data point \((\vec{x}, \vec{y}) \in \R^{n} \times \R^{p}\), we evaluate the loss on this data point as \(\ell(\vec{f}(\vec{x}, \bm{\theta}), \vec{y})\). The true ``empirical loss'' across a batch of \(q\) data points \((\vec{x}_{i}, \vec{y}_{i})\) is 
\begin{equation}
    L(\bm{\theta}) \doteq \sum_{i = 1}^{q}\ell(\vec{f}(\vec{x}_{i}, \bm{\theta}), \vec{y}_{i}),
\end{equation}
and a ``training algorithm'' attempts to optimize \(\bm{\theta}\) to minimize \(L\).


For now, we concentrate on the case where we have \(q = 1\), i.e., we have a single data point \((\vec{x}, \vec{y})\). Here we have 
\begin{equation}
    L(\bm{\theta}) = \ell(\vec{f}(\vec{x}, \bm{\theta}), \vec{y}).
\end{equation}
We now explore how to compute \(\nabla L(\bm{\theta})\).  In this situation, we are only differentiating with respect to \(\bm{\theta}\), so we fix \(\vec{x}\) and \(\vec{y}\), which gives the following computational graph:\footnote{Note that we did not break the graph down to matrix multiplications, vector additions, and applications of \(\vec{\sigma}^{(i)}\); we do this for clarity, because the full graph would be rather messy.}
\begin{figure}
    \centering
    \begin{tikzpicture}
        \node (f0i) [draw, circle, minimum size=1cm] {\(f_{i}^{(0)}\)};
        \node (f01) [draw, circle, minimum size=1cm, above=of f0i] {\(f_{1}^{(0)}\)};
        \node (f0n) [draw, circle, minimum size=1cm, below=of f0i] {\(f_{n_{0}}^{(1)}\)};
        
        \node (f11) [draw, circle, minimum size=1cm, right=of f01] {\(f_{1}^{(1)}\)};
        \node (f1i) [draw, circle, minimum size=1cm, right=of f0i] {\(f_{j}^{(1)}\)};
        \node (f1n) [draw, circle, minimum size=1cm, right=of f0n] {\(f_{n_{1}}^{(1)}\)};
        
        \node (d1) [right=of f11] {\(\cdots\)};
        \node (di) [right=of f1i] {\(\cdots\)};
        \node (dn) [right=of f1n] {\(\cdots\)};
        
        \node (fm1) [draw, circle, minimum size=1cm, right=of d1] {\(f_{i}^{(m)}\)};
        \node (fmi) [draw, circle, minimum size=1cm, right=of di] {\(f_{k}^{(m)}\)};
        \node (fmn) [draw, circle, minimum size=1cm, right=of dn] {\(f_{n_{m}}^{(m)}\)};
        
        \node (loss) [right=of fmi] {\(\ell(f(\vec{x}, \bm{\theta}), \vec{y})\)};

        \node (theta0) [below=of f0n] {\((W^{(0)}, \vec{b}^{(0)})\)};
        \node (theta1) [below=of f1n] {\((W^{(1)}, \vec{b}^{(1)})\)};
        \node (thetam) [below=of fmn] {\((W^{(m)}, \vec{b}^{(m)})\)};

        \node (z0) at ($(f01) + (0, 1)$) {\(\vec{z}^{(0)}(\vec{x}, \bm{\theta})\)};
        \node (z1) at ($(f11) + (0, 1)$) {\(\vec{z}^{(1)}(\vec{x}, \bm{\theta})\)};
        \node (zm) at ($(fm1) + (0, 1)$) {\(\vec{z}^{(m)}(\vec{x}, \bm{\theta})\)};

        \foreach \l in {0, 1, m} {
            \foreach \i in {1, i, n} {
                \draw [->, color=blue, bend left] (theta\l) to (f\l\i);
            }
        }

        \foreach \i in {1, i, n} {
            \foreach \j in {1, i, n} {
                \draw (f0\i) -- (f1\j);
            }
        }

        \foreach \i in {1, i, n} {
            \foreach \j in {1, i, n} {
                \draw (f1\i) -- (d\j);
            }
        }

        \foreach \i in {1, i, n} {
            \foreach \j in {1, i, n} {
                \draw (d\i) -- (fm\j);
            }
        }

        \foreach \i in {1, i, n} {
            \draw (fm\i) -- (loss);
        }

        \foreach \l in {0, 1, m} {
            \node at ($(f\l1)!0.45!(f\l i)$) {\(\vdots\)};
            \node at ($(f\l i)!0.45!(f\l n)$) {\(\vdots\)};
        }
    \end{tikzpicture}
    \caption{Computational graph corresponding to the multi-layer perceptron.}
\end{figure}

Since we have not defined gradients with respect to matrices (yet), we will focus on computing \(\nabla_{\vec{b}^{(i)}}L(\bm{\theta})\), i.e., the gradient of the loss function \(L\) with respect to \(\vec{b}^{(i)}\), as well as \(\nabla_{\vec{z}^{(i)}}L(\bm{\theta})\), i.e., the gradient of the loss function \(L\) with respect to \(\vec{z}^{(i)}(\vec{x}, \bm{\theta})\). They may be interpreted as sub-vectors (i.e., literally subsets of the entries) of the full gradient \(\nabla L(\bm{\theta}) = \nabla_{\bm{\theta}}L(\bm{\theta})\). This subscript notation is common in more involved problems which demand derivatives with respect to multiple vector-valued variables.

We first compute \(\nabla_{\vec{z}^{(m)}}L(\bm{\theta})\). We have 
\begin{equation}
    \nabla_{\vec{z}^{(m)}}L(\bm{\theta}) = \nabla_{\vec{z}^{(m)}}\ell(\vec{z}^{(m)}, \vec{y}).
\end{equation}
Indeed we cannot simplify this further, it being just the gradient of \(\ell\) in its first argument.

Now we handle general \(i \in \{1, \dots, m\}\). By chain rule we have the recursion:
\begin{align}
    \nabla_{\vec{z}^{(i - 1)}}L(\bm{\theta})
    &= [D_{\vec{z}^{(i - 1)}}L(\bm{\theta})]^{\top} \\ 
    &= [\{D_{\vec{z}^{(i)}}L(\bm{\theta})\}\{D_{\vec{z}^{(i - 1)}}\vec{z}^{(i)}(\vec{x}, \bm{\theta})\}]^{\top} \\ 
    &= [D_{\vec{z}^{(i - 1)}}\vec{z}^{(i)}(\vec{x}, \bm{\theta})]^{\top}[\nabla_{\vec{z}^{(i)}}L(\bm{\theta})].
\end{align}
Now we want to compute \(D_{\vec{z}^{(i - 1)}}\vec{z}^{(i)}\). Recall that 
\begin{equation}
    \vec{z}^{(i)} = W^{(i)}\vec{\sigma}^{(i)}(\vec{z}^{(i - 1)}) + \vec{b}^{(i)}.
\end{equation}
For convenience, write 
\begin{equation}
    \vec{u}^{(i)} \doteq \vec{\sigma}^{(i)}(\vec{z}^{(i - 1)}), \qquad \text{so that} \quad \vec{z}^{(i)} = W^{(i)}\vec{u}^{(i)} + \vec{b}^{(i)}.
\end{equation}
Now we compute derivatives, obtaining by the chain rule:
\begin{align}
    D_{\vec{z}^{(i - 1)}}\vec{z}^{(i)}
    &= [D_{\vec{u}^{(i)}}\vec{z}^{(i)}][D_{\vec{z}^{(i - 1)}}\vec{u}^{(i)}] \\
    &= W^{(i)}\cdot D\vec{\sigma}^{(i)}(\vec{z}^{(i - 1)}).
\end{align}
This gives us 
\begin{equation}
    D_{\vec{z}^{(i - 1)}}\vec{z}^{(i)} = W^{(i)}\cdot D\vec{\sigma}^{(i)}(\vec{z}^{(i - 1)}),
\end{equation}
and so 
\begin{equation}
    \nabla_{\vec{z}^{(i - 1)}}L(\bm{\theta}) = [W^{(i)} \cdot D\vec{\sigma}^{(i)}(\vec{z}^{(i - 1)})]^{\top}[\nabla_{\vec{z}^{(i)}}L(\bm{\theta})].
\end{equation}

To see the true value of this computation, we now compute \(\nabla_{\vec{b}^{(i)}}L(\bm{\theta})\), for each \(i \in \{0, \dots, m\}\). By another application of chain rule, we obtain
\begin{align}
    \nabla_{\vec{b}^{(i)}}L(\bm{\theta})
    &= [D_{\vec{b}^{(i)}}L(\bm{\theta})]^{\top} \\
    &= [\{D_{\vec{z}^{(i)}}L(\bm{\theta})\}\{D_{\vec{b}^{(i)}}\vec{z}^{(i)}\}]^{\top} \\ 
    &= [\{D_{\vec{z}^{(i)}}L(\bm{\theta})\}\cdot I]^{\top} \\  
    &= [D_{\vec{z}^{(i)}}L(\bm{\theta})]^{\top} \\
    &= \nabla_{\vec{z}^{(i)}}L(\bm{\theta}).
\end{align}

Now to do automatic differentiation by backpropagation, your favorite machine learning framework (such as PyTorch) computes all the gradients \(\nabla_{\vec{z}^{(i)}}L(\bm{\theta})\), starting at \(i = m\) and recursing until \(i = 0\). \textit{Computing the gradients in this order saves a lot of work, since each derivative is only computed once --- this is the main idea of backpropagation}. Once these derivatives are computed, one can use them to compute \(\nabla_{\vec{b}^{(i)}}L(\bm{\theta})\). If we can also compute \(\nabla_{W^{(i)}}L(\bm{\theta})\), then we can compute \(\nabla_{\bm{\theta}}L(\bm{\theta})\), and thus will be able to optimize \(L\) via gradient-based optimization algorithms such as gradient descent, as discussed later in the class. However, this computation will have to wait until slightly later when we cover matrix calculus.

We promised in this example a way to compute \(\nabla_{\bm{\theta}}L(\bm{\theta})\), or more precisely a way to compute \(\nabla_{W^{(i)}}L(\bm{\theta})\). We now have the tools to do this using the chain rule. Recall that we have access to \(\nabla_{\vec{z}^{(i)}}L(\bm{\theta})\) by backpropagation. Then we can compute the components of \(\nabla_{W^{(i)}}L(\bm{\theta})\) by
\begin{align}
    \nabla_{W^{(i)}}L(\bm{\theta})]_{j, k} 
    &= \pdv{L}{(W^{(i)})_{j, k}}(\bm{\theta}) \\
    &= \sum_{a}[\nabla_{\vec{z}^{(i)}}L(\bm{\theta})]_{a}\pdv{(\vec{z}^{(i)})_{a}}{(W^{(i)})_{j, k}} \\ 
    &= \sum_{a}[\nabla_{\vec{z}^{(i)}}L(\bm{\theta})]_{a}\cdot\begin{cases}[\vec{\sigma}^{(i)}(\vec{z}^{(i - 1)})]_{k}, & \text{if}\ a = j\ \text{and}\ i \in \{1, \dots, m\} \\ x_{k}, & \text{if}\ a = j\ \text{and}\ i = 0 \\ 0, & \text{otherwise}\end{cases} \\ 
    &= \begin{cases}[\nabla_{\vec{z}^{(i)}}L(\bm{\theta})]_{j}[\vec{\sigma}^{(i)}(\vec{z}^{(i - 1)})]_{k}, & \text{if}\ i \in \{1, \dots, m\} \\ [\nabla_{\vec{z}^{(i)}}L(\bm{\theta})]_{j}[\vec{x}]_{k}, & \text{if}\ i = 0\end{cases} \\ 
    &= \begin{cases}[\{\nabla_{\vec{z}^{(i)}}L(\bm{\theta})\}\vec{\sigma}^{(i)}(\vec{z}^{(i - 1)})^{\top}]_{j, k}, & \text{if}\ i \in \{1, \dots, m\} \\ \{\nabla_{\vec{z}^{(i)}}L(\bm{\theta})\}\vec{x}^{\top}]_{j, k}, & \text{if}\ i = 0.\end{cases}
\end{align}
This gives 
\begin{equation}
    \nabla_{W^{(i)}}L(\bm{\theta}) = \begin{cases}[\nabla_{\vec{z}^{(i)}}L(\bm{\theta})]\vec{\sigma}^{(i)}(\vec{z}^{(i - 1)})^{\top}, & \text{if}\ i \in \{1, \dots, m\} \\ [\nabla_{\vec{z}^{(i)}}L(\bm{\theta})]\vec{x}^{\top}, & \text{if}\ i = 0.\end{cases}
\end{equation}
In combination with the expression for \(\nabla_{\vec{b}^{(i)}}\), we can efficiently compute \(\nabla_{\bm{\theta}}L(\bm{\theta})\), and are able to train our neural network via gradient-based optimization methods such as gradient descent.




\section{Minmax Optimization}


\end{document}